% v1 draft, written by claude-sonnet-5-5 on 2026-11-12 for the SBSI 2027 Industry and Innovation
% track (TII-SI). Section 1 is drawn from the vault's debate notes; sections 2 and 3 are
% placeholders to be written from the product work. Every \fillin is a gap the authors must close.
%
% Template: the call requires the SBC article format. Download modelosparapublicaodeartigos.zip
% from https://www.sbc.org.br/ and put sbc-template.sty (and sbc.bst) next to this file. Until
% then the file falls back to plain article so it still compiles.
\documentclass[12pt]{article}

\usepackage[utf8]{inputenc}
\usepackage[T1]{fontenc}
\usepackage[english]{babel}
\usepackage{xcolor}
\usepackage{url}
\IfFileExists{sbc-template.sty}{\usepackage{sbc-template}}{\usepackage[margin=2.5cm]{geometry}}

% Marks a gap in the draft. Remove every use before submission.
\newcommand{\fillin}[1]{\textcolor{red}{\textbf{[TO FILL: #1]}}}

\title{[PRODUCT NAME]: Automatic Vulnerability Detection and Compliance Checking for AI-Generated Code}

\author{[AUTHOR 1], [AUTHOR 2]\\
[AFFILIATION]\\
\texttt{[emails]}}

\begin{document}

\maketitle

\begin{abstract}
AI coding assistants and agents now write a growing share of production code, and studies that measure the result report that a large fraction of it is vulnerable, in estimates that range from a substantial minority to most functionally correct output, depending on the benchmark and detector. Off-the-shelf static analysis is reported to miss many of these vulnerabilities, and prompting for security gives gains that vary by setup. This pitch presents [PRODUCT NAME], a \fillin{one-line description of the tool} that \fillin{what it detects and checks, and against which standards}. We built \fillin{prototype or proof of concept} and measured it on \fillin{benchmark or case}. We report \fillin{headline result and benefit for adopters}.
\end{abstract}

\section{Context}

AI assistants and agents are now routinely used to write code, and a developer who accepts a solution once it passes functional tests may leave vulnerabilities in the repository without noticing. The literature that measures this does not give a single number, because studies count different things with different detectors.

The first study of this kind, on the August 2021 GitHub Copilot preview, found that approximately 40\% of the programs it generated were vulnerable, across 89 scenarios and 1,689 programs~\cite{pearce2021asleep}. That figure divides by valid completions and does not assess functional correctness. Later work counts only functionally correct output. BaxBench, which runs expert-written exploits against 11 LLMs' backends, reports that on average exploits succeed against around half of the correct programs each model generates, and that the best correct-and-secure rate is about 37\% against a best pass@1 of 62\%~\cite{vero2025baxbench}. On 186 repository-level tasks, Zhao et al.\ report that for SWE-Agent with Claude 4 Sonnet 57.0\% of tasks are solved functionally but only 11.8\% are also secure, so 79.3\% of its functionally correct solutions are insecure~\cite{zhao2026vibe}. These estimates are not directly comparable: they cover different generators, tasks, denominators and ways of deciding that code is vulnerable.

Detection is where the practical problem sits. In BaxBench, Snyk Code missed 64 vulnerable programs that the benchmark's exploits showed to be vulnerable, among the correct programs generated by o1 (n = 237)~\cite{vero2025baxbench}. In RealSec-bench, CodeQL alerts reviewed by the authors reached a precision of 44.9\% on 89 alerts, and a two-stage voter-plus-judge adjudication lifted this to 81.7\% at a reported recall of 98.0\%~\cite{wang2026realsec}. The sources use different references for what counts as a missed vulnerability, so these results indicate the range of behaviour of existing detectors rather than a ranking of them.

Asking the model for secure code does not remove the need to check it. Zhao et al.\ report that, in one setting, adding a CWE hint raised the secure share from 11.8\% to 14.5\% (self-selected CWEs) and 15.1\% (oracle CWE) while functional pass fell from 57.0\% to 50.0\% and 55.4\%~\cite{zhao2026vibe}. A poster by Ma et al.\ reports that security-aware prompting reduced vulnerability rates by 37\% to up to 45\% across models, without stating whether the reductions are relative or absolute~\cite{ma2026rethinking}. The interventions and measures differ between these studies.

Within the 18 sources we reviewed, none evaluates compliance checking of AI-generated code against secure-coding standards such as CWE and OWASP, which appear only as taxonomies for choosing scenarios and exploits. \fillin{one sentence stating the problem this pitch addresses and the audience, e.g.\ teams adopting AI coding assistants who need verification that is cheap enough to run on every change}

\section{Process}

\fillin{how [PRODUCT NAME] was designed and tested; the call asks for the process adopted, for example proof of concept, prototype, measurement of indicators}

\begin{itemize}
  \item \fillin{what the tool does and how: detection techniques used (static analysis, LLM-based review, fuzzing, agentic red teaming), and how compliance with standards is checked}
  \item \fillin{what it was run on: benchmark, dataset or real case, with size and period}
  \item \fillin{baselines it was compared against, and the indicators measured (for example recall, precision, time per scan)}
  \item \fillin{who used it and in what setting, if any practitioner or company took part}
\end{itemize}

\section{Solution}

\fillin{results and benefits for the organisation and the software market; the jury criteria are clarity of description, degree or potential of innovation, and potential interest to companies, startups and entrepreneurs, so state the benefit for adopters directly}

\begin{itemize}
  \item \fillin{headline results, with numbers and the setting they hold in}
  \item \fillin{what is innovative compared with the detectors above}
  \item \fillin{limits of the evidence and next steps}
  \item \fillin{optional: link to one of the Grand Challenges of IS Research in Brazil 2026--2036, only if one genuinely fits}
\end{itemize}

\bibliographystyle{plain} % switch to sbc once sbc.bst is available
\bibliography{references}

\end{document}
